% =====================================================================
%  CV — example for a FICTIONAL person (Alex Sample).
%  Compile: pdflatex cv.tex   (run twice, for the page counter)
%
%  HOW TO TAILOR THIS FILE
%   1. Set \cvtitle to the exact job title in the posting.
%   2. Set \cvcountry to the posting's country (required — the build
%      fails without it; it drives the work-eligibility line).
%   3. Rewrite the PROFILE paragraph for that role (2-3 lines).
%   4. Reorder CORE COMPETENCIES so the posting's must-haves appear in
%      the first two groups; reorder KEY STRENGTHS the same way.
%   5. Lead PROFESSIONAL EXPERIENCE with the most relevant role.
%   6. Target two pages (~550-600 extracted words). Going long spills
%      onto page 3; cut content rather than shrinking the type.
%   7. Photo is off by default. \cvphoto{headshot.jpg} turns it on.
% =====================================================================
\documentclass[12pt,a4paper]{article}
\usepackage{jobcv}

% ------------------------------------------------------------ identity
\cvname{Alex Sample, M.Sc.}
\cvtitle{Embedded Systems Engineer}
\cvlocation{Utrecht, NL}
\cvphone{+31 6 0000 0000}
\cvemail{alex.sample@example.com}
\cvsite{alexsample.example}{https://alexsample.example}
\cvlinkedin{linkedin.com/in/alex-sample}{https://www.linkedin.com/in/alex-sample/}
% \cvphoto{headshot.jpg}

\cvcitizenship{Portuguese}
\cveligiblein{Netherlands, Portugal, Germany, Belgium, Spain, France}
\cvcountry{Germany}          % country of THIS posting — build fails if unset

\begin{document}
\pagestyle{cvrest}\thispagestyle{cvfirst}
\makecvheader

\cvsection{Profile}
Embedded systems engineer with six years of firmware and board bring-up for battery-powered
sensor products, from first prototype to certified volume production. Comfortable owning a
product's electronics end to end and working across mechanical, software and test teams.

\cvsection{Key Strengths}
\compactcols{3}{%
\quality{End-to-end product ownership}
\quality{Low-power design}
\quality{Clear documentation}
\quality{Calm under deadlines}
\quality{Cross-team collaborator}
\quality{Data-driven debugging}
}

\cvsection{Core Competencies}
\skillgroup{Firmware}{C, C++17, FreeRTOS, Zephyr; STM32, nRF52, ESP32; bootloaders, OTA updates;
SPI, I\textsuperscript{2}C, UART, BLE, LoRaWAN.}

\skillgroup{Hardware}{Schematic capture and PCB layout in KiCad and Altium Designer; power budgets,
battery management, EMC pre-compliance.}

\skillgroup{Test \& automation}{Python test fixtures, pytest, hardware-in-the-loop rigs; oscilloscope,
logic analyser, power profiler; CI with GitHub Actions.}

\skillgroup{Process}{Design reviews, DFM with contract manufacturers, CE/RED certification support,
Git, Jira.}

\cvsection{Professional Experience}
\cventry{Senior Embedded Engineer}{Example Sensors B.V., Utrecht, NL}{2022 -- present}
\begin{cvbullets}
  \item Owned firmware and electronics for a LoRaWAN soil-moisture sensor shipped to 40,000 sites.
  \item Cut average current draw from 48~µA to 11~µA, extending battery life from 2 to 7 years.
  \item Built a Python hardware-in-the-loop rig running 300 regression tests on every merge.
\end{cvbullets}

\cventry{Embedded Software Engineer}{Fictional Devices Lda., Porto, PT}{2019 -- 2022}
\begin{cvbullets}
  \item Developed BLE firmware for a wearable heart-rate monitor through CE certification.
\end{cvbullets}

\cventry{Graduate Engineer}{Placeholder Automation, Lisbon, PT}{2018 -- 2019}
\begin{cvbullets}
  \item Wrote PLC and HMI logic for packaging lines; commissioned three lines on site.
\end{cvbullets}

\cvsection{Education}
\eduentrywide{M.Sc., Electrical and Computer Engineering}{University of Example, Porto}{2016 -- 2018}
\eduentrywide{B.Sc., Electrical Engineering}{University of Example, Porto}{2013 -- 2016}

\cvsection{Languages}
\langline{{\bfseries English} --- Fluent \textperiodcentered\ {\bfseries Portuguese} --- Native
\textperiodcentered\ {\bfseries Dutch} --- B1}

\cvsection{At a Glance}
\compactcols{3}{%
\quality{2 products to volume production}
\quality{1 conference paper}
\quality{\permitline}
}

\end{document}
